\section{Motivating Example}
\label{sec:example}

The \texttt{Triple.emptyArray()} method from the \texttt{commons-lang} library is shown at the top of Figure~\ref{triple:emptyArray}. 
From the documentation, a developer would expect the method to provide two main behaviours: First, the description and the method name itself make it clear that the method should return an empty array. Second, the \texttt{@return} block reinforces that the method should always return the same array instance, since it is defined as a singleton.

\begin{figure}[h]
\begin{lstlisting}[language=Java]
/** Returns the empty array singleton that can   
 * be assigned without compiler warning.
 * 
 * @param <L> the left element type.
 * @param <M> the middle element type.
 * @param <R> the right element type.
 * @return the empty array singleton that can
 *  be assigned without compiler warning.
 * @since 3.10
 */
public static Triple<L, M, R>[] emptyArray() {
  return (Triple<L, M, R>[]) EMPTY_ARRAY;
}

// Only developer-written test for emptyArray
@Test
void testEmptyArrayGenerics() {
  Triple<Integer, String, Boolean>[] empty =
    Triple.emptyArray();
  assertEquals(0, empty.length);
}
\end{lstlisting}
\caption{Documentation, implementation, and developer-written test for \texttt{emptyArray()} from the \texttt{commons-lang} project.}
\label{triple:emptyArray}
\end{figure}

The only developer-written test case for \texttt{emptyArray} is shown at the bottom of Figure~\ref{triple:emptyArray}. While this test case achieves 100\% coverage for the \texttt{emptyArray} implementation, it only validates the first expected behaviour: that the returned array is actually empty.
The existing test suite exhibits a \emph{gap} between the expected behaviours and existing suite: the suite fails to validate that the returned array is a singleton. 
% no existing test case validate the second expected behaviour, that the implementation actually returns the same empty array every time, as it is a singleton.
As a result, if a regression were introduced that changed the method to return a new empty array each time, this change would go undetected and could introduce new faults to users of \texttt{emptyArray()} who relied on the method returning a singleton instance of the empty array.

% \begin{figure}[h]
% \begin{lstlisting}[language=Java]
%     @Test
%     void testEmptyArrayGenerics() {
%         final Triple<Integer, String, Boolean>[] empty = Triple.emptyArray();
%         assertEquals(0, empty.length);
%     }
% \end{lstlisting}    
% \caption{Developer-written test case for \texttt{emptyArray}. This test case achieves 100\% line coverage for the method, but only validates that the returned array is empty. Whether the returned array is a singleton is not checked.
% }
% \label{triple:emptyArrayTest}
% \end{figure}

%The associated test suite for this library includes one test case for this method, shown in~\autoref{triple:emptyArray}. This test verifies that the returned array has length zero. Although this test executes the method and confirms that the returned array is empty, it does not check the documented behaviour of the array to be a singleton instance. 

This paper seeks to address the gap between the source code and test case above.
While the developer-written test suite appears to be relatively complete from a coverage point-of-view, the test cases do not actually validate all of the behaviours a developer using this method might expect they could rely upon.
%
Our approach, called \toolname, takes as input the source code and documentation for a method, along with its existing test suite, and identifies these behavioural gaps between existing tests and expected behaviours.
%\toolname then generates tests for those behaviours, along with a description of a gap the newly-generated test aims to validate.
These gaps are actionable: developers can resolve them by strengthening their test suite to directly address the specifically-identified weakness in their existing tests.
This gap can be easily addressed by a single \texttt{assertSame(..)} assertion across two invocations to validate that the returned array is the same instance for multiple calls.

%Figure~\ref{triple:emptyArrayGen} shows how the identified gap could be easily resolved by the developer with a simple additional test case for \texttt{emptyArray()} by validating that the returned array is the same instance every time.


%Our approach addresses this gap. \toolname analyzes the documentation and extracts two behaviours: (1) the returned array is empty, and (2) the returned array is a singleton. It then maps these documented behaviours to the existing test cases. In this example, the existing test already covers the first behaviour (emptiness), but no test verifies the second behaviour (singleton). As a result, \toolname identifies the singleton behaviour as untested and generates a targeted test, shown in~\autoref{triple:emptyArray_generated_test}, that explicitly checks this case. This generated test does not increase structural coverage, since the lines were already executed by the existing test. Instead, it strengthens the suite by ensuring that both documented behaviours are explicitly verified. The existing test confirms that the returned array is empty, while the generated test adds coverage for the singleton behaviour. Together, they ensure that the complete documented behaviour of this method is now tested.

% \begin{figure}[h]
% \begin{lstlisting}[language=Java]
% @Test
% void testEmptyArrayReturnsSingletonInstance() {
%   Triple<Integer, String, Boolean>[] array1 =
%     Triple.emptyArray();
%   Triple<Integer, String, Boolean>[] array2 =
%     Triple.emptyArray();
%   assertSame(array1, array2);
% }
% \end{lstlisting}
% \caption{An exemplar test case that resolves the \toolname{}-identified behavioural gap for \texttt{emptyArray()}.
% }
% \label{triple:emptyArrayGen}
% \end{figure}

The remainder of this paper describes an approach for detecting behavioural gaps, evaluates the accuracy of this technique, and describes the kinds of behaviours that are frequently missed by existing developer-written test suites.

%This example illustrates the need for documentation-driven testing. It shows how structural coverage can obscure untested behaviours, how developers lack straightforward ways to detect these semantic gaps, and how our approach can generate focused tests that bridge them.
